\documentclass[11pt,english]{article}
\pagestyle{plain}
\pagenumbering{arabic}
\usepackage[english]{babel}
\usepackage{graphicx}
\usepackage{multirow}
\usepackage{subfig}
%\usepackage{auto-pst-pdf}
\usepackage[table,xcdraw]{xcolor}
\usepackage[T1]{fontenc}
\usepackage{hyperref}
\usepackage[hmargin={0.8cm,1.2cm},vmargin=1.2cm]{geometry}
\usepackage{lscape}
\usepackage{pdflscape}
\renewcommand*\familydefault{\sfdefault}
\newcommand{\HRule}{\rule{\linewidth}{0.5mm}}
% packages for headers and footers
\usepackage{fancyhdr}
%\usepackage{lastpage}
\pagestyle{fancy}
\fancyhf{}
%\fancyfoot[C]{Page\ \thepage\ of \pageref{LastPage}}

\begin{document}

\begin{figure}
\begin{center}
\hspace{1.5cm}
\parbox{5.5cm}{\includegraphics[width=2.2cm]{/home/report/SCRIPTS/LOGOS/logoLaMMA}}
\hspace{.3cm}
\parbox{5.5cm}{\includegraphics[width=2.2cm]{/home/report/SCRIPTS/LOGOS/fate_logo_11def}}
\hspace{.3cm}
\parbox{5.5cm}{\includegraphics[width=2.2cm]{/home/report/SCRIPTS/LOGOS/INAF_Arcetri_colore}}
\hspace{.1cm}
\vspace{1.2cm}
\end{center}
\end{figure}

\begin{center}
\HRule \\[0.4cm]
%\Huge{\textbf{Progetto FATE LaMMA-OAA per ESO @ Cerro Paranal etcetcetc\ldots}}
\Huge{\textbf{Forecast system for the turbulence and meteorology at the Paranal Observatory}}
\\[0.4cm]
\LARGE{\textbf{Performance assessment - monthly report}}
\HRule \\[0.4cm]
\end{center}

\begin{center}
\vspace{2cm}\Huge{Issued: Mon, 26-06-2023\\ 17:39 Paranal local time}
\HRule \\[0.1cm]
\end{center}

%\begin{center}
%\HRule \\[0.1cm]
%\tableofcontents
%\HRule \\[0.1cm]
%\end{center}
%\newpage
\clearpage

%%%%%%%%%%%%%%%%%%%%%%%%%%%%%%%%%
\begin{figure}
\centering
\subfloat[]{\includegraphics[width=.33\linewidth,angle=-90]{/home/report/FIGS/ws_sim_mnh_ar_dimm_300_960_BEF_stan.eps}}
\subfloat[]{\includegraphics[width=.33\linewidth,angle=-90]{/home/report/FIGS/ws_sim_mnh_ar_dimm_300_960_AFT_stan.eps}}
\subfloat[]{\includegraphics[width=.33\linewidth,angle=-90]{/home/report/FIGS/ws_sim_mnh_ar_dimm_300_960_AFT_stan_per.eps}}
\caption{{wind speed} ($m s^{-1}$): (a) STANDARD CONFIGURATION, (b) WITH AR (1H), (c) PERSISTENCE.}
\label{fig:ws}
\end{figure}
%%%%%%%%%%%%%%%%%%%%%%%%%%%%%%%%%
\begin{figure}
\centering
\subfloat[]{\includegraphics[width=.33\linewidth,angle=-90]{/home/report/FIGS/wd_sim_mnh_ar_dimm_300_960_BEF_stan.eps}}
\subfloat[]{\includegraphics[width=.33\linewidth,angle=-90]{/home/report/FIGS/wd_sim_mnh_ar_dimm_300_960_AFT_stan.eps}}
\subfloat[]{\includegraphics[width=.33\linewidth,angle=-90]{/home/report/FIGS/wd_sim_mnh_ar_dimm_300_960_AFT_stan_per.eps}}
\caption{wind direction (degree): (a) STANDARD CONFIGURATION, (b) WITH AR, (c) PERSISTENCE.}
\label{fig:wd}
\end{figure}
%%%%%%%%%%%%%%%%%%%%%%%%%%%%%%%%%
\begin{figure}
\centering
\subfloat[]{\includegraphics[width=.33\linewidth,angle=-90]{/home/report/FIGS/rh_sim_mnh_ar_dimm_300_960_BEF_stan.eps}}
\subfloat[]{\includegraphics[width=.33\linewidth,angle=-90]{/home/report/FIGS/rh_sim_mnh_ar_dimm_300_960_AFT_stan.eps}}
\subfloat[]{\includegraphics[width=.33\linewidth,angle=-90]{/home/report/FIGS/rh_sim_mnh_ar_dimm_300_960_AFT_stan_per.eps}}
\caption{relative humidity (percent): (a) STANDARD CONFIGURATION, (b) WITH AR (1H), (c) PERSISTENCE.}
\label{fig:rh}
\end{figure}
%%%%%%%%%%%%%%%%%%%%%%%%%%%%%%%%%
\begin{figure}
\centering
\subfloat[]{\includegraphics[width=.33\linewidth,angle=-90]{/home/report/FIGS/pwv_sim_mnh_ar_dimm_300_960_BEF_stan.eps}}
\subfloat[]{\includegraphics[width=.33\linewidth,angle=-90]{/home/report/FIGS/pwv_sim_mnh_ar_dimm_300_960_AFT_stan.eps}}
\subfloat[]{\includegraphics[width=.33\linewidth,angle=-90]{/home/report/FIGS/pwv_sim_mnh_ar_dimm_300_960_AFT_stan_per.eps}}
\caption{precipitable water vapor (mm): (a) STANDARD CONFIGURATION, (b) WITH AR (1H), (c) PERSISTENCE.}
\label{fig:pwv}
\end{figure}
%%%%%%%%%%%%%%%%%%%%%%%%%%%%%%%%%
\begin{figure}
\centering
\subfloat[]{\includegraphics[width=.33\linewidth,angle=-90]{/home/report/FIGS/see_sim_mnh_ar_dimm_300_960_BEF_os18_1000.eps}}
\subfloat[]{\includegraphics[width=.33\linewidth,angle=-90]{/home/report/FIGS/see_sim_mnh_ar_dimm_300_960_AFT_os18_1000.eps}}
\subfloat[]{\includegraphics[width=.33\linewidth,angle=-90]{/home/report/FIGS/see_sim_mnh_ar_dimm_300_960_AFT_os18_1000_per.eps}}
\caption{total seeing (arcsec): (a) STANDARD CONFIGURATION ($<$ 1.5''), (b) WITH AR (1H), (c) PERSISTENCE.}
\label{fig:see}
\end{figure}
%%%%%%%%%%%%%%%%%%%%%%%%%%%%%%%%%
\begin{figure}
\centering
\subfloat[]{\includegraphics[width=.33\linewidth,angle=-90]{/home/report/FIGS/tau_sim_mnh_ar_dimm_300_960_BEF_os18_1000.eps}}
\subfloat[]{\includegraphics[width=.33\linewidth,angle=-90]{/home/report/FIGS/tau_sim_mnh_ar_dimm_300_960_AFT_os18_1000.eps}}
\subfloat[]{\includegraphics[width=.33\linewidth,angle=-90]{/home/report/FIGS/tau_sim_mnh_ar_dimm_300_960_AFT_os18_1000_per.eps}}
\caption{coeherence time (ms): (a) STANDARD CONFIGURATION, (b) WITH AR (1H), (c) PERSISTENCE.}
\label{fig:tau}
\end{figure}
%%%%%%%%%%%%%%%%%%%%%%%%%%%%%%%%%
\begin{figure}
\centering
\subfloat[]{\includegraphics[width=.33\linewidth,angle=-90]{/home/report/FIGS/glf_sim_mnh_ar_dimm_300_960_BEF_os18.eps}}
\subfloat[]{\includegraphics[width=.33\linewidth,angle=-90]{/home/report/FIGS/glf_sim_mnh_ar_dimm_300_960_AFT_os18.eps}}
\subfloat[]{\includegraphics[width=.33\linewidth,angle=-90]{/home/report/FIGS/glf_sim_mnh_ar_dimm_300_960_AFT_os18_per.eps}}
\caption{ground layer fraction (dimensionless): (a) STANDARD CONFIGURATION, (b) WITH AR (1H), (c) PERSISTENCE.}
\label{fig:glf}
\end{figure}

%%%%%%%%%%%%%%%%%%%%%%%%%%%%%%%%%
\newpage
%%%%%%%%%%%%%%%%%%%%%%%%%%%%%%%%%

%%%%%%%%%%%%%%%%%%%%
\clearpage
%%%%%%%%%%%%%%%%%%%%
%%%%%%%%%%%%%%%%%%%%%%%%%%%%%%%%%%
%\newpage
\begin{table}[]
\begin{center}
\begin{tabular}{|l|l|l|l|l|l|l|}
\hline
\multicolumn{1}{|c|}{\cellcolor[HTML]{C0C0C0}\textbf{VARIABLE}} & \multicolumn{1}{c|}{\cellcolor[HTML]{C0C0C0}\textbf{BIAS}} & \multicolumn{1}{c|}{\cellcolor[HTML]{C0C0C0}\textbf{RMSE}} & \multicolumn{1}{c|}{\cellcolor[HTML]{C0C0C0}\textbf{SD}} & \multicolumn{1}{c|}{\cellcolor[HTML]{C0C0C0}\textbf{BIAS}} & \multicolumn{1}{c|}{\cellcolor[HTML]{C0C0C0}\textbf{RMSE}} & \multicolumn{1}{c|}{\cellcolor[HTML]{C0C0C0}\textbf{SD}}\\
\multicolumn{1}{|c|}{\cellcolor[HTML]{C0C0C0}} & \multicolumn{1}{c|}{\cellcolor[HTML]{C0C0C0}(inc. month)} & \multicolumn{1}{c|}{\cellcolor[HTML]{C0C0C0}(inc. month)} & \multicolumn{1}{c|}{\cellcolor[HTML]{C0C0C0}(inc. month)} & \multicolumn{1}{c|}{\cellcolor[HTML]{C0C0C0}(last month)} & \multicolumn{1}{c|}{\cellcolor[HTML]{C0C0C0}(last month)} & \multicolumn{1}{c|}{\cellcolor[HTML]{C0C0C0}(last month)}\\\hline
\cellcolor[HTML]{C0C0C0}WS (m$\cdot$s$^{-1}$) &      0.10  &      2.49  &      2.49  &     -0.61     &      2.60     &      2.53\\
\cellcolor[HTML]{C0C0C0}WD (degree)           &     -0.97  &     35.22  &     35.21  &      1.64     &     33.85     &     33.81\\
\cellcolor[HTML]{C0C0C0}RH (percent)          &      0.73  &      6.64  &      6.60  &      0.76     &      6.84     &      6.79\\
\cellcolor[HTML]{C0C0C0}PWV (mm)              &     -0.20 &      1.00 &      0.99 &     -0.35    &      1.00    &      0.94\\
\cellcolor[HTML]{C0C0C0}SEE (arcsec)          &     -0.02 &      0.23 &      0.23 &     -0.02    &      0.22    &      0.21\\
\cellcolor[HTML]{C0C0C0}TAU (ms)              &     -0.06 &      2.19 &      2.19 &     -0.63    &      2.32    &      2.24\\
\cellcolor[HTML]{C0C0C0}GLF (dimensionless)   &     -0.03 &      0.19 &      0.18 &     -0.01    &      0.16    &      0.16\\
\hline
\end{tabular}
\caption{Statistics for variables in standard configuration}
\end{center}
\end{table}

%%%%%%%%%%%%%%%%%%%%%%%%%%%%%%%%%%
%\newpage
\begin{table}[]
\begin{center}
\begin{tabular}{|l|l|l|l|l|l|l|}
\hline
\multicolumn{1}{|c|}{\cellcolor[HTML]{C0C0C0}\textbf{VARIABLE}} & \multicolumn{1}{c|}{\cellcolor[HTML]{C0C0C0}\textbf{BIAS}} & \multicolumn{1}{c|}{\cellcolor[HTML]{C0C0C0}\textbf{RMSE}} & \multicolumn{1}{c|}{\cellcolor[HTML]{C0C0C0}\textbf{SD}} & \multicolumn{1}{c|}{\cellcolor[HTML]{C0C0C0}\textbf{BIAS}} & \multicolumn{1}{c|}{\cellcolor[HTML]{C0C0C0}\textbf{RMSE}} & \multicolumn{1}{c|}{\cellcolor[HTML]{C0C0C0}\textbf{SD}}\\
\multicolumn{1}{|c|}{\cellcolor[HTML]{C0C0C0}} & \multicolumn{1}{c|}{\cellcolor[HTML]{C0C0C0}(inc. month)} & \multicolumn{1}{c|}{\cellcolor[HTML]{C0C0C0}(inc. month)} & \multicolumn{1}{c|}{\cellcolor[HTML]{C0C0C0}(inc. month)} & \multicolumn{1}{c|}{\cellcolor[HTML]{C0C0C0}(persistence)} & \multicolumn{1}{c|}{\cellcolor[HTML]{C0C0C0}(persistence)} & \multicolumn{1}{c|}{\cellcolor[HTML]{C0C0C0}(persistence)}\\\hline
\cellcolor[HTML]{C0C0C0}WS (m$\cdot$s$^{-1}$) &     -0.01  &      0.66  &      0.66  &     -0.02  &      0.92  &      0.92\\
\cellcolor[HTML]{C0C0C0}WD (degree)           &      0.30  &      6.65  &      6.64  &      1.90  &     12.47  &     12.33\\
\cellcolor[HTML]{C0C0C0}RH (percent)          &     -0.01  &      1.37  &      1.37  &     -0.05  &      2.18  &      2.18\\
\cellcolor[HTML]{C0C0C0}PWV (mm)              &     -0.01 &      0.13 &      0.13 &      0.04 &      0.25 &      0.25\\
\cellcolor[HTML]{C0C0C0}SEE (arcsec)          &     -0.00 &      0.10 &      0.10 &     -0.00 &      0.15 &      0.15\\
\cellcolor[HTML]{C0C0C0}TAU (ms)              &     -0.01 &      0.58 &      0.58 &     -0.03 &      1.06 &      1.06\\
\cellcolor[HTML]{C0C0C0}GLF (dimensionless)   &     -0.00 &      0.06 &      0.06 &     -0.01 &      0.09 &      0.09\\
\hline
\end{tabular}
\caption{Statistics for variables processed with AR (1H)}
\end{center}
\end{table}


%%%%%%%%%%%%%%%%%%%%%%%%%%%%%%%%%
\newpage
%%%%%%%%%%%%%%%%%%%%%%%%%%%%%%%%%

%%%%%%%%%%%%%%%%%%%%
\clearpage
%%%%%%%%%%%%%%%%%%%%
%%%%%%%%%%%%%%%%%%%%%%%%%%%%%%%%%
\begin{table}[]
\begin{center}
\begin{tabular}{llccc}
\hline
{wind speed}                                       &                                                    & \multicolumn{3}{c}{Observations}                 \\
{$m s^{-1}$}                                       &                             & ws $<$ 4.2   & 4.2 $<$ ws $<$ 7.5 & ws $>$ 7.5 \\
\hline
\multicolumn{1}{c}{\multirow{3}{*}{Model data}}  & ws $<$ 4.2             & 4541                & 2127                       & 249              \\
                                                 & 4.2  $<$ ws $<$ 7.5 & 3034                & 3970                       & 1347              \\
                                                 & ws $>$ 7.5             & 430                & 1908                       & 6408              \\
\hline 
\multicolumn{5}{l}{Sample size: 38880; PC=62.1\%; EBD=2.8\%; POD1=56.7\%; POD2=49.6\%; POD3=80.1\%}                 \\
\end{tabular}
\end{center}
\caption{Contingency table for {wind speed} ($m s^{-1}$) in standard configuration}
\label{tab:contingencywsBEF}
\end{table}
%%%%%%%%%%%%%%%%%%%%%%%%%%%%%%%%%
\begin{table}[]
\begin{center}
\begin{tabular}{llccc}
\hline
{wind speed}                                       &                                                    & \multicolumn{3}{c}{Observations}                 \\
{$m s^{-1}$}                                       &                             & ws $<$ 4.2   & 4.2 $<$ ws $<$ 7.5 & ws $>$ 7.5 \\
\hline
\multicolumn{1}{c}{\multirow{3}{*}{Model data}}  & ws $<$ 4.2             & 7377                & 623                       & 1              \\
                                                 & 4.2  $<$ ws $<$ 7.5 & 628                & 6950                       & 386              \\
                                                 & ws $>$ 7.5             & 0                & 432                       & 7617              \\
\hline 
\multicolumn{5}{l}{Sample size: 38880; PC=91.4\%; EBD=0.0\%; POD1=92.2\%; POD2=86.8\%; POD3=95.2\%}                 \\
\end{tabular}
\end{center}
\caption{Contingency table for variable wind speed ($m s^{-1}$) processed with AR (1H)}
\label{tab:contingencywsAFT}
\end{table}
%%%%%%%%%%%%%%%%%%%%%%%%%%%%%%%%%
\begin{table}[]
\begin{center}
\begin{tabular}{llcccc}
\hline
{wind direction}                                       &                             &    \multicolumn{4}{c}{Observations}                 \\
{degree}                                       &                             & {North}   & {East} & {South} & {West} \\
\hline
\multicolumn{1}{c}{\multirow{4}{*}{Model data}}  & {North}                    & 9401   & 1028    & 107     & 229\\
                                                 & {East}                     & 1002   & 1832    & 560     & 0\\
                                                 & {South}                    & 99   & 508   & 3567    & 195\\
                                                 & {West}                     & 202  & 89   & 63    & 199\\
\hline 
\multicolumn{6}{l}{Sample size: 38880; PC=78.6\%; EBD=1.5\%; POD1=87.8\%; POD2=53.0\%; POD3=83.0\%; POD4=31.9\%}                 \\
\end{tabular}
\end{center}
\caption{Contingency table for {wind direction} (degree) in standard configuration}
\label{tab:contingencywdBEF}
\end{table}
%%%%%%%%%%%%%%%%%%%%%%%%%%%%%%%%%
\begin{table}[]
\begin{center}
\begin{tabular}{llcccc}
\hline
{wind direction}                                       &                                                    & \multicolumn{4}{c}{Observations}                 \\
{degree}                                       &                             & {North}   & {East} & {South} & {West} \\
\hline
\multicolumn{1}{c}{\multirow{4}{*}{Model data}}  & {North}              & 10453                & 222                       & 0              & 36\\
                                                 & {East}               & 217                & 3151                       & 40              & 0\\
                                                 & {South}              & 0                & 84                      & 4236             & 29\\
                                                 & {West}               & 34               & 0                      & 21             & 558\\
\hline 
\multicolumn{6}{l}{Sample size: 38880; PC=96.4\%; EBD=0.0\%; POD1=97.7\%; POD2=91.1\%; POD3=98.6\%; POD4=89.6\%}                 \\
\end{tabular}
\end{center}
\caption{Contingency table for {wind direction} (degree) in standard configuration}
\label{tab:contingencywdAFT}
\end{table}

%%%%%%%%%%%%%%%%%%%%%%%%%%%%%%%%%
\newpage
%%%%%%%%%%%%%%%%%%%%%%%%%%%%%%%%%

%%%%%%%%%%%%%%%%%%%%
\clearpage
%%%%%%%%%%%%%%%%%%%%
%%%%%%%%%%%%%%%%%%%%%%%%%%%%%%%%%
\begin{table}[]
\begin{center}
\begin{tabular}{llccc}
\hline
{relative humidity}                                       &                                                    & \multicolumn{3}{c}{Observations}                 \\
{percent}                                       &                             & rh $<$ 7.1   & 7.1 $<$ rh $<$ 15.3 & rh $>$ 15.3 \\
\hline
\multicolumn{1}{c}{\multirow{3}{*}{Model data}}  & rh $<$ 7.1             & 4755                & 1095                       & 27              \\
                                                 & 7.1  $<$ rh $<$ 15.3 & 2696                & 5193                       & 1357              \\
                                                 & rh $>$ 15.3             & 379                & 1539                       & 6443              \\
\hline
\multicolumn{5}{l}{Sample size: 38232; PC=69.8\%; EBD=1.7\%; POD1=60.7\%; POD2=66.3\%; POD3=82.3\%}
\end{tabular}
\end{center}
\caption{Contingency table for variable relative humidity (percent) in standard configuration}
\label{tab:contingencyrhBEF}
\end{table}
%%%%%%%%%%%%%%%%%%%%%%%%%%%%%%%%%
\begin{table}[]
\begin{center}
\begin{tabular}{llccc}
\hline
{relative humidity}                                       &                                                    & \multicolumn{3}{c}{Observations}                 \\
{percent}                                       &                             & rh $<$ 7.1   & 7.1 $<$ rh $<$ 15.3 & rh $>$ 15.3 \\
\hline
\multicolumn{1}{c}{\multirow{3}{*}{Model data}}  & rh $<$ 7.1             & 7408                & 328                       & 1              \\
                                                 & 7.1  $<$ rh $<$ 15.3 & 422                & 7294                       & 248              \\
                                                 & rh $>$ 15.3             & 0                & 205                       & 7578              \\
\hline
\multicolumn{5}{l}{Sample size: 38232; PC=94.9\%; EBD=0.0\%; POD1=94.6\%; POD2=93.2\%; POD3=96.8\%}
\end{tabular}
\end{center}
\caption{Contingency table for variable relative humidity (percent) processed with AR (1H)}
\label{tab:contingencyrhAFT}
\end{table}
%%%%%%%%%%%%%%%%%%%%%%%%%%%%%%%%%
\begin{table}[]
\begin{center}
\begin{tabular}{llccc}
\hline
{precipitable water vapor}                                       &                                                    & \multicolumn{3}{c}{Observations}                 \\
{mm}                                       &                             & pwv $<$ 1.9   & 1.9 $<$ pwv $<$ 3.5 & pwv $>$ 3.5 \\
\hline
\multicolumn{1}{c}{\multirow{3}{*}{Model data}}  & pwv $<$ 1.9             & 3402                & 854                       & 2              \\
                                                 & 1.9  $<$ pwv $<$ 3.5 & 443                & 2596                       & 576              \\
                                                 & pwv $>$ 3.5             & 5                & 400                       & 3271              \\
\hline
\multicolumn{5}{l}{Sample size: 18900; PC=80.3\%; EBD=0.1\%; POD1=88.4\%; POD2=67.4\%; POD3=85.0\%}                 \\
\end{tabular}
\end{center}
\caption{Contingency table for variable precipitable water vapor (mm) in standard configuration}
\label{tab:contingencypwvBEF}
\end{table}
%%%%%%%%%%%%%%%%%%%%%%%%%%%%%%%%%
\begin{table}[]
\begin{center}
\begin{tabular}{llccc}
\hline
{precipitable water vapor}                                       &                                                    & \multicolumn{3}{c}{Observations}                 \\
{mm}                                       &                             & pwv $<$ 1.9   & 1.9 $<$ pwv $<$ 3.5 & pwv $>$ 3.5 \\
\hline
\multicolumn{1}{c}{\multirow{3}{*}{Model data}}  & pwv $<$ 1.9             & 3785                & 69                       & 0              \\
                                                 & 1.9  $<$ pwv $<$ 3.5 & 65                & 3733                       & 39              \\
                                                 & pwv $>$ 3.5             & 0                & 48                       & 3810              \\
\hline 
\multicolumn{5}{l}{Sample size: 18900; PC=98.1\%; EBD=0.0\%; POD1=98.3\%; POD2=97.0\%; POD3=99.0\%}
\end{tabular}
\end{center}
\caption{Contingency table for variable precipitable water vapor (mm) processed with AR (1H)}
\label{tab:contingencypwvAFT}
\end{table}

%%%%%%%%%%%%%%%%%%%%%%%%%%%%%%%%%
\newpage
%%%%%%%%%%%%%%%%%%%%%%%%%%%%%%%%%

%%%%%%%%%%%%%%%%%%%%
\clearpage
%%%%%%%%%%%%%%%%%%%%
%%%%%%%%%%%%%%%%%%%%%%%%%%%%%%%%%
\begin{table}[]
\begin{center}
\begin{tabular}{llccc}
\hline
{total seeing}                                       &                                                    & \multicolumn{3}{c}{Observations}                 \\
{arcsec}                                       &                             & see $<$ 0.6   & 0.6 $<$ see $<$ 0.8 & see $>$ 0.8 \\
\hline
\multicolumn{1}{c}{\multirow{3}{*}{Model data}}  & see $<$ 0.6             & 3678                & 2340                       & 841              \\
                                                 & 0.6  $<$ see $<$ 0.8 & 1828                & 2382                       & 1886              \\
                                                 & see $>$ 0.8             & 520                & 1303                       & 3298              \\
\hline
\multicolumn{5}{l}{Sample size: 37584; PC=51.8\%; EBD=7.5\%; POD1=61.0\%; POD2=39.5\%; POD3=54.7\%}
\end{tabular}
\end{center}
\caption{Contingency table for variable total seeing (arcsec) in standard configuration and accuracy 0.0}
\label{tab:contingencyseeBEF}
\end{table}
%%%%%%%%%%%%%%%%%%%%%%%%%%%%%%%%%
\begin{table}[]
\begin{center}
\begin{tabular}{llccc}
\hline
{total seeing}                                       &                                                    & \multicolumn{3}{c}{Observations}                 \\
{arcsec}                                       &                             & see $<$ 0.6   & 0.6 $<$ see $<$ 0.8 & see $>$ 0.8 \\
\hline
\multicolumn{1}{c}{\multirow{3}{*}{Model data}}  & see $<$ 0.6             & 5603                & 0                       & 597              \\
                                                 & 0.6  $<$ see $<$ 0.8 & 0                & 5649                       & 597              \\
                                                 & see $>$ 0.8             & 423                & 376                       & 5428              \\
\hline
\multicolumn{5}{l}{Sample size: 37584; PC=92.3\%; EBD=5.6\%; POD1=93.0\%; POD2=93.8\%; POD3=82.0\%}
\end{tabular}
\end{center}
\caption{Contingency table for variable total seeing (arcsec) in standard configuration and accuracy 0.24}
\label{tab:contingencyseeBEF}
\end{table}
%%%%%%%%%%%%%%%%%%%%%%%%%%%%%%%%%
\begin{table}[]
\begin{center}
\begin{tabular}{llccc}
\hline
{total seeing}                                       &                                                    & \multicolumn{3}{c}{Observations}                 \\
{arcsec}                                       &                             & see $<$ 0.6   & 0.6 $<$ see $<$ 0.8 & see $>$ 0.8 \\
\hline
\multicolumn{1}{c}{\multirow{3}{*}{Model data}}  & see $<$ 0.6             & 5707                & 835                       & 23              \\
                                                 & 0.6  $<$ see $<$ 0.8 & 725                & 4942                       & 546              \\
                                                 & see $>$ 0.8             & 14                & 669                       & 5877              \\
\hline
\multicolumn{5}{l}{Sample size: 37584; PC=85.5\%; EBD=0.2\%; POD1=88.5\%; POD2=76.7\%; POD3=91.2\%}
\end{tabular}
\end{center}
\caption{Contingency table for variable total seeing (arcsec) processed with AR (1H) and accuracy 0.0}
\label{tab:contingencyseeAFT}
\end{table}
%%%%%%%%%%%%%%%%%%%%%%%%%%%%%%%%%
\begin{table}[]
\begin{center}
\begin{tabular}{llccc}
\hline
{total seeing}                                       &                                                    & \multicolumn{3}{c}{Observations}                 \\
{arcsec}                                       &                             & see $<$ 0.6   & 0.6 $<$ see $<$ 0.8 & see $>$ 0.8 \\
\hline
\multicolumn{1}{c}{\multirow{3}{*}{Model data}}  & see $<$ 0.6             & 6433                & 11                       & 23              \\
                                                 & 0.6  $<$ see $<$ 0.8 & 0                & 6404                       & 18              \\
                                                 & see $>$ 0.8             & 13                & 31                       & 6423              \\
\hline
\multicolumn{5}{l}{Sample size: 37584; PC=99.6\%; EBD=0.2\%; POD1=99.8\%; POD2=99.3\%; POD3=99.4\%}
\end{tabular}
\end{center}
\caption{Contingency table for variable total seeing (arcsec) processed with AR (1H) and accuracy 0.24}
\label{tab:contingencyseeAFT}
\end{table}

%%%%%%%%%%%%%%%%%%%%%%%%%%%%%%%%%
\newpage
%%%%%%%%%%%%%%%%%%%%%%%%%%%%%%%%%

%%%%%%%%%%%%%%%%%%%%
\clearpage
%%%%%%%%%%%%%%%%%%%%
%%%%%%%%%%%%%%%%%%%%%%%%%%%%%%%%%
\begin{table}[]
\begin{center}
\begin{tabular}{llccc}
\hline
{coeherence time}                                       &                                                    & \multicolumn{3}{c}{Observations}                 \\
{ms}                                       &                             & tau $<$ 3.4   & 3.4 $<$ tau $<$ 5.8 & tau $>$ 5.8 \\
\hline
\multicolumn{1}{c}{\multirow{3}{*}{Model data}}  & tau $<$ 3.4             & 4336                & 1463                       & 107              \\
                                                 & 3.4  $<$ tau $<$ 5.8 & 1221                & 3060                       & 1597              \\
                                                 & tau $>$ 5.8             & 196                & 1230                       & 4049              \\
\hline
\multicolumn{5}{l}{Sample size: 37476; PC=66.3\%; EBD=1.8\%; POD1=75.4\%; POD2=53.2\%; POD3=70.4\%}
\end{tabular}
\end{center}
\caption{Contingency table for variable coeherence time (ms) in standard configuration and accuracy 0.0}
\label{tab:contingencytauBEF}
\end{table}
%%%%%%%%%%%%%%%%%%%%%%%%%%%%%%%%%
\begin{table}[]
\begin{center}
\begin{tabular}{llccc}
\hline
{coeherence time}                                       &                                                    & \multicolumn{3}{c}{Observations}                 \\
{ms}                                       &                             & tau $<$ 3.4   & 3.4 $<$ tau $<$ 5.8 & tau $>$ 5.8 \\
\hline
\multicolumn{1}{c}{\multirow{3}{*}{Model data}}  & tau $<$ 3.4             & 5229                & 149                       & 639              \\
                                                 & 3.4  $<$ tau $<$ 5.8 & 328                & 4937                       & 635              \\
                                                 & tau $>$ 5.8             & 196                & 667                       & 5114              \\
\hline
\multicolumn{5}{l}{Sample size: 37476; PC=88.5\%; EBD=4.8\%; POD1=90.9\%; POD2=85.8\%; POD3=80.1\%}
\end{tabular}
\end{center}
\caption{Contingency table for variable coeherence time (ms) in standard configuration and accuracy 1.22}
\label{tab:contingencytauBEF}
\end{table}
%%%%%%%%%%%%%%%%%%%%%%%%%%%%%%%%%
\begin{table}[]
\begin{center}
\begin{tabular}{llccc}
\hline
{coeherence time}                                       &                                                    & \multicolumn{3}{c}{Observations}                 \\
{ms}                                       &                             & tau $<$ 3.4   & 3.4 $<$ tau $<$ 5.8 & tau $>$ 5.8 \\
\hline
\multicolumn{1}{c}{\multirow{3}{*}{Model data}}  & tau $<$ 3.4             & 5386                & 339                       & 2              \\
                                                 & 3.4  $<$ tau $<$ 5.8 & 366                & 5011                       & 400              \\
                                                 & tau $>$ 5.8             & 1                & 403                       & 5351              \\
\hline
\multicolumn{5}{l}{Sample size: 37476; PC=91.2\%; EBD=0.0\%; POD1=93.6\%; POD2=87.1\%; POD3=93.0\%}
\end{tabular}
\end{center}
\caption{Contingency table for variable coeherence time (ms) processed with AR (1H) and accuracy 0.0}
\label{tab:contingencytauAFT}
\end{table}
%%%%%%%%%%%%%%%%%%%%%%%%%%%%%%%%%
\begin{table}[]
\begin{center}
\begin{tabular}{llccc}
\hline
{coeherence time}                                       &                                                    & \multicolumn{3}{c}{Observations}                 \\
{ms}                                       &                             & tau $<$ 3.4   & 3.4 $<$ tau $<$ 5.8 & tau $>$ 5.8 \\
\hline
\multicolumn{1}{c}{\multirow{3}{*}{Model data}}  & tau $<$ 3.4             & 5745                & 0                       & 33              \\
                                                 & 3.4  $<$ tau $<$ 5.8 & 7                & 5725                       & 33              \\
                                                 & tau $>$ 5.8             & 1                & 28                       & 5720              \\
\hline
\multicolumn{5}{l}{Sample size: 37476; PC=99.6\%; EBD=0.2\%; POD1=99.9\%; POD2=99.5\%; POD3=98.9\%}
\end{tabular}
\end{center}
\caption{Contingency table for variable coeherence time (ms) processed with AR (1H) and accuracy 1.22}
\label{tab:contingencytauAFT}
\end{table}

%%%%%%%%%%%%%%%%%%%%%%%%%%%%%%%%%
\newpage
%%%%%%%%%%%%%%%%%%%%%%%%%%%%%%%%%

%%%%%%%%%%%%%%%%%%%%
\clearpage
%%%%%%%%%%%%%%%%%%%%
%%%%%%%%%%%%%%%%%%%%%%%%%%%%%%%%%
\begin{table}[]
\begin{center}
\begin{tabular}{llccc}
\hline
{ground layer fraction}                                       &                                                    & \multicolumn{3}{c}{Observations}                 \\
{dimensionless}                                       &                             & glf $<$ 0.6   & 0.6 $<$ glf $<$ 0.7 & glf $>$ 0.7 \\
\hline
\multicolumn{1}{c}{\multirow{3}{*}{Model data}}  & glf $<$ 0.6             & 3194                & 2361                       & 1636              \\
                                                 & 0.6  $<$ glf $<$ 0.7 & 1856                & 2275                       & 2456              \\
                                                 & glf $>$ 0.7             & 707                & 1121                       & 1665              \\
\hline
\multicolumn{5}{l}{Sample size: 37476; PC=41.3\%; EBD=13.6\%; POD1=55.5\%; POD2=39.5\%; POD3=28.9\%}
\end{tabular}
\end{center}
\caption{Contingency table for variable ground layer fraction (dimensionless) in standard configuration and accuracy 0.0}
\label{tab:contingencyglfBEF}
\end{table}
%%%%%%%%%%%%%%%%%%%%%%%%%%%%%%%%%
\begin{table}[]
\begin{center}
\begin{tabular}{llccc}
\hline
{ground layer fraction}                                       &                                                    & \multicolumn{3}{c}{Observations}                 \\
{dimensionless}                                       &                             & glf $<$ 0.6   & 0.6 $<$ glf $<$ 0.7 & glf $>$ 0.7 \\
\hline
\multicolumn{1}{c}{\multirow{3}{*}{Model data}}  & glf $<$ 0.6             & 4964                & 566                       & 1758              \\
                                                 & 0.6  $<$ glf $<$ 0.7 & 86                & 5104                       & 1493              \\
                                                 & glf $>$ 0.7             & 707                & 87                       & 3999              \\
\hline
\multicolumn{5}{l}{Sample size: 37476; PC=81.4\%; EBD=14.3\%; POD1=86.2\%; POD2=88.7\%; POD3=55.2\%}
\end{tabular}
\end{center}
\caption{Contingency table for variable ground layer fraction (dimensionless) in standard configuration and accuracy 0.14}
\label{tab:contingencyglfBEF}
\end{table}
%%%%%%%%%%%%%%%%%%%%%%%%%%%%%%%%%
\begin{table}[]
\begin{center}
\begin{tabular}{llccc}
\hline
{ground layer fraction}                                       &                                                    & \multicolumn{3}{c}{Observations}                 \\
{dimensionless}                                       &                             & glf $<$ 0.6   & 0.6 $<$ glf $<$ 0.7 & glf $>$ 0.7 \\
\hline
\multicolumn{1}{c}{\multirow{3}{*}{Model data}}  & glf $<$ 0.6             & 4991                & 837                       & 27              \\
                                                 & 0.6  $<$ glf $<$ 0.7 & 726                & 4180                       & 842              \\
                                                 & glf $>$ 0.7             & 40                & 740                       & 4888              \\
\hline
\multicolumn{5}{l}{Sample size: 37476; PC=81.4\%; EBD=0.4\%; POD1=86.7\%; POD2=72.6\%; POD3=84.9\%}
\end{tabular}
\end{center}
\caption{Contingency table for variable ground layer fraction (dimensionless) processed with AR (1H) and accuracy 0.0}
\label{tab:contingencyglfAFT}
\end{table}
%%%%%%%%%%%%%%%%%%%%%%%%%%%%%%%%%
\begin{table}[]
\begin{center}
\begin{tabular}{llccc}
\hline
{ground layer fraction}                                       &                                                    & \multicolumn{3}{c}{Observations}                 \\
{dimensionless}                                       &                             & glf $<$ 0.6   & 0.6 $<$ glf $<$ 0.7 & glf $>$ 0.7 \\
\hline
\multicolumn{1}{c}{\multirow{3}{*}{Model data}}  & glf $<$ 0.6             & 5707                & 42                       & 36              \\
                                                 & 0.6  $<$ glf $<$ 0.7 & 10                & 5694                       & 35              \\
                                                 & glf $>$ 0.7             & 40                & 21                       & 5721              \\
\hline
\multicolumn{5}{l}{Sample size: 37476; PC=99.1\%; EBD=0.4\%; POD1=99.1\%; POD2=98.9\%; POD3=98.8\%}
\end{tabular}
\end{center}
\caption{Contingency table for variable ground layer fraction (dimensionless) processed with AR (1H) and accuracy 0.14}
\label{tab:contingencyglfAFT}
\end{table}

%%%%%%%%%%%%%%%%%%%%%%%%%%%%%%%%%
\newpage
%%%%%%%%%%%%%%%%%%%%%%%%%%%%%%%%%

%%%%%%%%%%%%%%%%%%%%
\clearpage
%%%%%%%%%%%%%%%%%%%%
%%%%%%%%%%%%%%%%%%%%%%%%%%%%%%%%%%
\begin{table}[]
\begin{center}
\begin{tabular}{|l|l|l|l|}
\hline
\multicolumn{1}{|c|}{\cellcolor[HTML]{C0C0C0}\textbf{PARAMETER}} & \multicolumn{1}{c|}{\cellcolor[HTML]{C0C0C0}\textbf{STANDARD}} & \multicolumn{1}{c|}{\cellcolor[HTML]{C0C0C0}\textbf{WITH AR (1H)}} \\
\hline
\cellcolor[HTML]{C0C0C0}POD1  & 56.7                                & 92.2         \\
\cellcolor[HTML]{C0C0C0}POD2  & 49.6                                & 86.8         \\
\cellcolor[HTML]{C0C0C0}POD3  & 80.1                                & 95.2         \\
\cellcolor[HTML]{C0C0C0}PC    & 62.1                                  & 91.4           \\
\cellcolor[HTML]{C0C0C0}EBD   & 2.8                                 & 0.0          \\
\hline
\end{tabular}
\caption{PODs for wind speed ($m s^{-1}$)}
\end{center}
\end{table}
%%%%%%%%%%%%%%%%%%%%%%%%%%%%%%%%%%
\begin{table}[]
\begin{center}
\begin{tabular}{|l|l|l|l|}
\hline
\multicolumn{1}{|c|}{\cellcolor[HTML]{C0C0C0}\textbf{PARAMETER}} & \multicolumn{1}{c|}{\cellcolor[HTML]{C0C0C0}\textbf{STANDARD}} & \multicolumn{1}{c|}{\cellcolor[HTML]{C0C0C0}\textbf{WITH AR (1H)}} \\
\hline
\cellcolor[HTML]{C0C0C0}POD1  & 87.8      & 97.7     \\
\cellcolor[HTML]{C0C0C0}POD2  & 53.0      & 91.1     \\
\cellcolor[HTML]{C0C0C0}POD3  & 83.0      & 98.6     \\
\cellcolor[HTML]{C0C0C0}POD4  & 31.9      & 89.6     \\
\cellcolor[HTML]{C0C0C0}PC    & 78.6        & 96.4       \\
\cellcolor[HTML]{C0C0C0}EBD   & 1.5       & 0.0      \\
\hline
\end{tabular}
\caption{PODs for wind direction (degree)}
\end{center}
\end{table}
%%%%%%%%%%%%%%%%%%%%%%%%%%%%%%%%%%
\begin{table}[]
\begin{center}
\begin{tabular}{|l|l|l|l|}
\hline
\multicolumn{1}{|c|}{\cellcolor[HTML]{C0C0C0}\textbf{PARAMETER}} & \multicolumn{1}{c|}{\cellcolor[HTML]{C0C0C0}\textbf{STANDARD}} & \multicolumn{1}{c|}{\cellcolor[HTML]{C0C0C0}\textbf{WITH AR (1H)}} \\
\hline
\cellcolor[HTML]{C0C0C0}POD1  & 60.7                                & 94.6         \\
\cellcolor[HTML]{C0C0C0}POD2  & 66.3                                & 93.2         \\
\cellcolor[HTML]{C0C0C0}POD3  & 82.3                                & 96.8         \\
\cellcolor[HTML]{C0C0C0}PC    & 69.8                                  & 94.9           \\
\cellcolor[HTML]{C0C0C0}EBD   & 1.7                                 & 0.0          \\
\hline
\end{tabular}
\caption{PODs for relative humidity (percent)}
\end{center}
\end{table}
%%%%%%%%%%%%%%%%%%%%%%%%%%%%%%%%%%
\begin{table}[]
\begin{center}
\begin{tabular}{|l|l|l|l|}
\hline
\multicolumn{1}{|c|}{\cellcolor[HTML]{C0C0C0}\textbf{PARAMETER}} & \multicolumn{1}{c|}{\cellcolor[HTML]{C0C0C0}\textbf{STANDARD}} & \multicolumn{1}{c|}{\cellcolor[HTML]{C0C0C0}\textbf{WITH AR (1H)}} \\
\hline
\cellcolor[HTML]{C0C0C0}POD1  & 88.4                                & 98.3         \\
\cellcolor[HTML]{C0C0C0}POD2  & 67.4                                & 97.0         \\
\cellcolor[HTML]{C0C0C0}POD3  & 85.0                                & 99.0         \\
\cellcolor[HTML]{C0C0C0}PC    & 80.3                                  & 98.1           \\
\cellcolor[HTML]{C0C0C0}EBD   & 0.1                                 & 0.0          \\
\hline
\end{tabular}
\caption{PODs for precipitable water vapor (mm)}
\end{center}
\end{table}

%%%%%%%%%%%%%%%%%%%%%%%%%%%%%%%%%
\newpage
%%%%%%%%%%%%%%%%%%%%%%%%%%%%%%%%%

%%%%%%%%%%%%%%%%%%%%
\clearpage
%%%%%%%%%%%%%%%%%%%%
%%%%%%%%%%%%%%%%%%%%%%%%%%%%%%%%%%
\begin{table}[]
\begin{center}
\begin{tabular}{|l|l|l|l|}
\hline
\multicolumn{1}{|c|}{\cellcolor[HTML]{C0C0C0}\textbf{PARAMETER}} & \multicolumn{1}{c|}{\cellcolor[HTML]{C0C0C0}\textbf{STANDARD}} & \multicolumn{1}{c|}{\cellcolor[HTML]{C0C0C0}\textbf{WITH AR (1H)}} \\
\hline
\cellcolor[HTML]{C0C0C0}POD1  & 61.0                                & 88.5         \\
\cellcolor[HTML]{C0C0C0}POD2  & 39.5                                & 76.7         \\
\cellcolor[HTML]{C0C0C0}POD3  & 54.7                                & 91.2         \\
\cellcolor[HTML]{C0C0C0}PC    & 51.8                                  & 85.5           \\
\cellcolor[HTML]{C0C0C0}EBD   & 7.5                                 & 0.2          \\
\hline
\end{tabular}
\caption{PODs for total seeing (arcsec) and accuray 0.0}
\end{center}
\end{table}
%%%%%%%%%%%%%%%%%%%%%%%%%%%%%%%%%%
\begin{table}[]
\begin{center}
\begin{tabular}{|l|l|l|l|}
\hline
\multicolumn{1}{|c|}{\cellcolor[HTML]{C0C0C0}\textbf{PARAMETER}} & \multicolumn{1}{c|}{\cellcolor[HTML]{C0C0C0}\textbf{STANDARD}} & \multicolumn{1}{c|}{\cellcolor[HTML]{C0C0C0}\textbf{WITH AR (1H)}} \\
\hline
\cellcolor[HTML]{C0C0C0}POD1  & 93.0                                & 99.8         \\
\cellcolor[HTML]{C0C0C0}POD2  & 93.8                                & 99.3         \\
\cellcolor[HTML]{C0C0C0}POD3  & 82.0                                & 99.4         \\
\cellcolor[HTML]{C0C0C0}PC    & 92.3                                  & 99.6           \\
\cellcolor[HTML]{C0C0C0}EBD   & 5.6                                 & 0.2          \\
\hline
\end{tabular}
\caption{PODs for total seeing (arcsec) and accuray 0.24}
\end{center}
\end{table}
%%%%%%%%%%%%%%%%%%%%%%%%%%%%%%%%%%
\begin{table}[]
\begin{center}
\begin{tabular}{|l|l|l|l|}
\hline
\multicolumn{1}{|c|}{\cellcolor[HTML]{C0C0C0}\textbf{PARAMETER}} & \multicolumn{1}{c|}{\cellcolor[HTML]{C0C0C0}\textbf{STANDARD}} & \multicolumn{1}{c|}{\cellcolor[HTML]{C0C0C0}\textbf{WITH AR (1H)}} \\
\hline
\cellcolor[HTML]{C0C0C0}POD1  & 75.4                                & 93.6         \\
\cellcolor[HTML]{C0C0C0}POD2  & 53.2                                & 87.1         \\
\cellcolor[HTML]{C0C0C0}POD3  & 70.4                                & 93.0         \\
\cellcolor[HTML]{C0C0C0}PC    & 66.3                                  & 91.2           \\
\cellcolor[HTML]{C0C0C0}EBD   & 1.8                                 & 0.0          \\
\hline
\end{tabular}
\caption{PODs for coeherence time (ms) and accuray 0.0}
\end{center}
\end{table}
%%%%%%%%%%%%%%%%%%%%%%%%%%%%%%%%%%
\begin{table}[]
\begin{center}
\begin{tabular}{|l|l|l|l|}
\hline
\multicolumn{1}{|c|}{\cellcolor[HTML]{C0C0C0}\textbf{PARAMETER}} & \multicolumn{1}{c|}{\cellcolor[HTML]{C0C0C0}\textbf{STANDARD}} & \multicolumn{1}{c|}{\cellcolor[HTML]{C0C0C0}\textbf{WITH AR (1H)}} \\
\hline
\cellcolor[HTML]{C0C0C0}POD1  & 90.9                                & 99.9         \\
\cellcolor[HTML]{C0C0C0}POD2  & 85.8                                & 99.5         \\
\cellcolor[HTML]{C0C0C0}POD3  & 80.1                                & 98.9         \\
\cellcolor[HTML]{C0C0C0}PC    & 88.5                                  & 99.6           \\
\cellcolor[HTML]{C0C0C0}EBD   & 4.8                                 & 0.2          \\
\hline
\end{tabular}
\caption{PODs for coeherence time (ms) and accuray 1.22}
\end{center}
\end{table}
%%%%%%%%%%%%%%%%%%%%%%%%%%%%%%%%%%
\begin{table}[]
\begin{center}
\begin{tabular}{|l|l|l|l|}
\hline
\multicolumn{1}{|c|}{\cellcolor[HTML]{C0C0C0}\textbf{PARAMETER}} & \multicolumn{1}{c|}{\cellcolor[HTML]{C0C0C0}\textbf{STANDARD}} & \multicolumn{1}{c|}{\cellcolor[HTML]{C0C0C0}\textbf{WITH AR (1H)}} \\
\hline
\cellcolor[HTML]{C0C0C0}POD1  & 55.5                                & 86.7         \\
\cellcolor[HTML]{C0C0C0}POD2  & 39.5                                & 72.6         \\
\cellcolor[HTML]{C0C0C0}POD3  & 28.9                                & 84.9         \\
\cellcolor[HTML]{C0C0C0}PC    & 41.3                                  & 81.4           \\
\cellcolor[HTML]{C0C0C0}EBD   & 13.6                                 & 0.4          \\
\hline
\end{tabular}
\caption{PODs for ground layer fraction (dimensionless) and accuray 0.0}
\end{center}
\end{table}
%%%%%%%%%%%%%%%%%%%%%%%%%%%%%%%%%%
\begin{table}[]
\begin{center}
\begin{tabular}{|l|l|l|l|}
\hline
\multicolumn{1}{|c|}{\cellcolor[HTML]{C0C0C0}\textbf{PARAMETER}} & \multicolumn{1}{c|}{\cellcolor[HTML]{C0C0C0}\textbf{STANDARD}} & \multicolumn{1}{c|}{\cellcolor[HTML]{C0C0C0}\textbf{WITH AR (1H)}} \\
\hline
\cellcolor[HTML]{C0C0C0}POD1  & 86.2                                & 99.1         \\
\cellcolor[HTML]{C0C0C0}POD2  & 88.7                                & 98.9         \\
\cellcolor[HTML]{C0C0C0}POD3  & 55.2                                & 98.8         \\
\cellcolor[HTML]{C0C0C0}PC    & 81.4                                  & 99.1           \\
\cellcolor[HTML]{C0C0C0}EBD   & 14.3                                 & 0.4          \\
\hline
\end{tabular}
\caption{PODs for ground layer fraction (dimensionless) and accuray 0.14}
\end{center}
\end{table}
%%%%%%%%%%%%%%%%%%%%
\clearpage
%%%%%%%%%%%%%%%%%%%%
\begin{center}
\HRule \\[0.4cm]
SHORT-TERM FORECAST
\HRule \\[0.4cm]
\end{center}

computed each full hour for each parameter

Night forecast parameters:
\begin{itemize}
        \item seeing (see)
        \item tau0 (tau)
        \item ground-layer fraction (glf)
        \item precipitable water vapor (pwv)
        \item relative humidity (rh)
        \item wind speed (ws)
        \item wind direction (wd)
\end{itemize}
Day forecast parameters:
\begin{itemize}
        \item precipitable water vapor (pwv)
        \item relative humidity (rh)
        \item wind speed (ws)
        \item wind direction (wd)      
\end{itemize}
This forecast relies on measurements obtained each full hour from http://archive.eso.org/wdb, using WDB interface
Forecast (for each parameter) is only available if the following conditions are met:
\begin{enumerate}
\item http://archive.eso.org/wdb is accessible and available
\item at least 2 past days worth of measurements available for the specific parameter
\item last measurement (at each full hour) not older than 10 minutes
\item at least 30 minutes of measurements in the current time span (sunset-sunrise for night forecast and sunrise-sunset for day forecast)
\end{enumerate}

FILE NAME FORMAT:
\begin{verbatim}
YYYY-MM-DDTHH_KIND_AR_output_TIME.csv
(All dates are in UT)
YYYY=year
MM=month
DD=day
HH=hour
\end{verbatim}
The above dates refer to the time at which the short-term forecast is computed and are the starting point of the forecast extended on 4 hours in the future

KIND=One of the following:
\begin{itemize}
\item    see\_evol\_time = seeing time evolution
\item    tau\_evol\_time = tau0 time evolution
\item    glf\_evol\_time = glf time evolution
\item    pwv\_evol\_time = pwv time evolution
\item    rh\_k6\_evol\_time = rh time evolution at 30m a.g.l.
\item    ws\_k6\_evol\_time = ws time evolution at 30m a.g.l.
\item    wd\_k6\_evol\_time = wd time evolution at 30m a.g.l.
\end{itemize}
TIME=One of the following:
\begin{itemize}
\item    night = night forecasts
\item    day = day forecast
\end{itemize}
 % FINO A QUI TUTTO BENE

FILE HEADER:
\begin{itemize}
\item        2 comment lines:
  \begin{itemize}
  \item                BUFFER LENGTH = X, DATA AGE(s) = Y --> X is the number of past days/nights considered for the computation of the forecast (minimum 2, maximum 5)
                                                       Performances for the AR in Masciadri et al. 2023 have been calculated for X=5.
                                                       Y is how old is the last measurement (in seconds) at COMPUTE TIME  !!!!!!!!   XXXXXXX  !!!!!!!
  \item                COMPUTE TIME = YYYY-MM-DDTHH:mm:ss  ---> time at which the forecast is computed up to minute (mm) and second (ss)
  \end{itemize}
\item        header line for the 2 columns (each file has 2 columns)
  \begin{itemize}
  \item                Forecast time (\%Y-\%m-\%dT\%H:\%M:\%S) = time for which the forecast refers to (extended up to 4 hours from COMPUTE TIME)
  \item                PARAMETER = the parameter which is forecasted (seeing, tau0, rh, etc...)
  \end{itemize}
\end{itemize}
%

\begin{verbatim}
EXAMPLE OUTPUTS:
2023-05-23T09_glf_evol_time_AR_output_night.csv
2023-05-23T09_pwv_evol_time_AR_output_night.csv
2023-05-23T09_rh_k6_evol_time_AR_output_night.csv
2023-05-23T09_see_evol_time_AR_output_night.csv
2023-05-23T09_tau_evol_time__AR_output_night.csv
2023-05-23T09_wd_k6_evol_time_AR_output_night.csv
2023-05-23T09_ws_k6_evol_time_AR_output_night.csv
\end{verbatim}
MEASUREMENTS REFERENCE:
\begin{itemize}
\item Seeing reference observations: DIMM
\item Tau and glf reference observations: MASS
\item Water vapor reference observations: LHATPRO (A)
\item Relative humidity, wind speed and direction are taken from the sensors at 30m
\end{itemize}
All obsrvations come from \href{http://archive.eso.org/cms/eso-data/ambient-conditions/paranal-ambient-query-forms.html}{the archive eso-data website}
%
PS: the convention used for the forecast is described in Masciadri et al. 2023-Session 3. When we compute for example the forecast at 1h it includes all values between 0h and 1h. If we consider the forecast at 2h we take all values between 1h and 2h. And so on.
%

%%%%%%%%%%%%%%%%%%%%%%%%%%%%%%
\begin{center}
\HRule \\[0.4cm]
VALIDATION TEST (NO AR=autoregression)
\HRule \\[0.4cm]
\end{center}

There are deterministic and probabilistic forecasts as requested by ESO. For all the parameters (with exception of the cloud cover) we treat the outputs as:
\begin{itemize}
\item deterministic OUTPUTS: one file collecting the forecasts of all the parameters 
\item probabilistic OUTPUTS: individual files for each parameter 
\end{itemize}

For the cloud cover the deterministic and probabilistic forecasts are both in the same file as this parameter has to be treated differently from the other parameters according to the SOW.

FILE NAME FORMAT:

\verb+AAAAAAAA-BB_CCCCCCCC-DDDD_EEEEEEEE-FFFF_TYPE_KIND_NIGHT.csv+
(All dates are in UT)
\begin{itemize}
\item AAAAAAAA= date in which the initialisation and forcing data for the mesoscale model are calculated (date format: YYYYMMDD)
\item BB=hour in which initialisation and forcing data are calculated (date format: HH)
\item CCCCCCCC=Date corresponding to the start of the simulated period i.e.to the sunset (date format: YYYYMMDD)
\item DDDD=Hour and minute associated to the date CCCCCCCC (hour format: HHMM)
\item EEEEEEEE=Date corresponding to the end of simulated period i.e. to the sunrise (date format: YYYYMMDD)
\item FFFF=Hour and minute associated to the date EEEEEEEE (hour format: HHMM)
\end{itemize}

TYPE=Kind of forecast ---> 
\begin{itemize}	
\item probLONG=probability that a forecast is included in specific ranges of values
\item detLONG=determinist forecasts
\item probdetLONG= probabilist and deterministic forecast for the cloud cover (the cloud cover has to be treated differently from the other parameters - see SOW)
\end{itemize}	
%
%
PS: The term LONG is associated to the forecast delivered at maximum 2 hours before the sunset (this is the time within which we have to provide the forecast on a time scale of 1, 2 and 3 nights according to the ESO documentation - Enclosure 2-2). These are the forecast requested for the Validation test.

%
KIND=One of the following:
%
\begin{itemize}	
\item cloud=deterministic and probabilistic forecast of the cloud cover that is treated differently from all the other parameters (see SOW).
\item standardALL=determinist forecast produced by Astro-Meso-NH in standard method related to *all* the parameters: see, tau0, glf, ws, wd, rh, pwv.
\item glf,pwv,rh,see,tau0see,wd,ws = for the probabilistic forecasts we have used the specific ranges/intervals/thresholds for each parameter indicated in SOW (September 2020 version):
\begin{itemize}
\item glf: section 4.2.3
\item pwv: section 4.2.7
\item rh: section 4.2.6:
\item see:section 4.2.1
\item tau0see: section 4.2.2
\item Wd: section 4.2.5
\item Ws: section 4.2.4
\end{itemize}
\end{itemize}
%       
%
%
NIGHT=One of the following:
\begin{itemize}
\item night1=relative to the forecast on a time scale of 1d in the future (timestep=10m)
\item night2=relative to the forecast on a time scale of 2d in the future (timestep=1h)
\item night3=relative to the forecast on a time scale of 3d in the future (timestep=3h)
\end{itemize}
%
%
EXAMPLE OUTPUTS (each night has 27 outputs -> 9 outputs x 3 nights):
\begin{verbatim}
1. 20210111-00_20210111-2339_20210112-1000_detLONG_standardALL_night1.csv (7 figures)
2. 20210111-00_20210111-2339_20210112-1000_probLONG_glf_night1.csv (1 figure)
3. 20210111-00_20210111-2339_20210112-1000_probLONG_pwv_night1.csv (1 figure)
4. 20210111-00_20210111-2339_20210112-1000_probLONG_rh_night1.csv  (1 figure)
5. 20210111-00_20210111-2339_20210112-1000_probLONG_see_night1.csv.(1 figure)
6. 20210111-00_20210111-2339_20210112-1000_probLONG_tau0see_night1.csv. (1 figure)
7. 20210111-00_20210111-2339_20210112-1000_probLONG_wd_night1.csv. (1 figure)
8. 20210111-00_20210111-2339_20210112-1000_probLONG_ws_night1.csv. (1 figure)
9. 20210111-00_20210111-2339_20210112-1000_probdetLONG_cloud_night1.csv (2 figures)
10. 20210111-00_20210112-2339_20210113-1001_detLONG_standardALL_night2.csv
11. 20210111-00_20210112-2339_20210113-1001_probLONG_glf_night2.csv
12. 20210111-00_20210112-2339_20210113-1001_probLONG_pwv_night2.csv
13. 20210111-00_20210112-2339_20210113-1001_probLONG_rh_night2.csv
14. 20210111-00_20210112-2339_20210113-1001_probLONG_see_night2.csv
15. 20210111-00_20210112-2339_20210113-1001_probLONG_tau0see_night2.csv
16. 20210111-00_20210112-2339_20210113-1001_probLONG_wd_night2.csv
17. 20210111-00_20210112-2339_20210113-1001_probLONG_ws_night2.csv
18. 20210111-00_20210112-2339_20210113-1001_probdetLONG_cloud_night2.csv
19. 20210111-00_20210113-2339_20210114-1001_detLONG_standardALL_night3.csv
20. 20210111-00_20210113-2339_20210114-1001_probLONG_glf_night3.csv
21. 20210111-00_20210113-2339_20210114-1001_probLONG_pwv_night3.csv
22. 20210111-00_20210113-2339_20210114-1001_probLONG_rh_night3.csv
23. 20210111-00_20210113-2339_20210114-1001_probLONG_see_night3.csv
24. 20210111-00_20210113-2339_20210114-1001_probLONG_tau0see_night3.csv
25. 20210111-00_20210113-2339_20210114-1001_probLONG_wd_night3.csv
26. 20210111-00_20210113-2339_20210114-1001_probLONG_ws_night3.csv
27. 20210111-00_20210113-2339_20210114-1001_probdetLONG_cloud_night3.csv
\end{verbatim}

FILES HEADER 

The headers are mostly auto-esplicative.
%
We add here some more infos:
%
\begin{itemize}
\item  The files reporting the probability that seeing and tau0 fall within specific categories are indicated with Category 1, 2, 3, 4, 5, 6 and 7 according to file readme\_fig4.png and according to the SOW (September update).
\item According to SOW, section 4.1.6, the cloud cover is treated differently than the other parameters. Files whose name is *probdet\_cloud* include the determinist and probabilistic forecast according to Section 4.1.6 of SOW and the digital document AD2.
\item According to SOW - section 4.2.6 we include the probability that the wind direction blows from N, S, W and E. Meteo Convention: 0 deg = North, 90 deg = East. In our analysis we divided the 360 degrees in 4 sectors: "North": means that the wind direction is included between -45 degrees and + 45 degrees. Same logic for the others cardinal directions. 
\item  The ranges for the calculation of the probability that all the atmospheric and astroclimatic parameters assume specific values match tables reported in SOW.
\end{itemize}
%%%%%


%%%%%%%%%%%%%%%%%%%%%%%%%%%%%%
\begin{center}
\HRule \\[0.4cm]
VALIDATION TEST (NO AR=autoregression)
\HRule \\[0.4cm]
\end{center}

There are deterministic and probabilistic forecasts as requested by ESO. For all the parameters (with exception of the cloud cover) we treat the outputs as:
\begin{itemize}
\item         deterministic OUTPUTS: one file collecting the forecasts of all the parameters 
             
\item probabilistic OUTPUTS: individual files for each parameter 
\end{itemize}


For the cloud cover the deterministic and probabilistic forecasts are both in the same file as this parameter has to be treated differently from the other parameters according to the SOW.


FILE NAME FORMAT:

\verb+AAAAAAAA-BB_CCCCCCCC-DDDD_EEEEEEEE-FFFF_TYPE_KIND_NIGHT.csv+
(All dates are in UT)
\begin{itemize}
\item AAAAAAAA= date in which the initialisation and forcing data for the mesoscale model are calculated (date format: YYYYMMDD)
\item BB=hour in which initialisation and forcing data are calculated (date format: HH)
\item CCCCCCCC=Date corresponding to the start of the simulated period i.e.to the sunrise (date format: YYYYMMDD)
\item DDDD=Hour and minute associated to the date CCCCCCCC (hour format: HHMM)
\item EEEEEEEE=Date corresponding to the end of simulated period i.e. to the sunset (date format: YYYYMMDD)
\item FFFF=Hour and minute associated to the date EEEEEEEE (hour format: HHMM)
\end{itemize}


TYPE=Kind of forecast ---> 
\begin{itemize}
\item probLONG=probability that a forecast is included in specific ranges of values
\item detLONG=determinist forecasts
\item probdetLONG= probabilist and deterministic forecast for the cloud cover (the cloud cover has to be treated differently from the other parameters - see SOW)
\end{itemize}


PS: The term LONG is associated to the forecast delivered at maximum 2 hours before the sunset (this is the time within which we have to provide the forecast on a time scale of 1, 2 and 3 nights according to the ESO documentation - Enclosure 2-2). These are the forecast requested for the Validation test.


KIND=One of the following:
\begin{itemize}
\item cloud=deterministic and probabilistic forecast of the cloud cover that is treated differently from all the other parameters (see SOW).
\item standardALL=determinist forecast produced by Astro-Meso-NH in standard method related to *all* the parameters: ws, wd, rh, pwv.
\item wd,ws = for the probabilistic forecasts we have used the specific ranges/intervals/thresholds for each parameter indicated in SOW (September 2020 version):
\begin{itemize}
\item pwv: section 4.2.7
\item rh: section 4.2.6:
\item Wd: section 4.2.5
\item Ws: section 4.2.4
\end{itemize}
\end{itemize}


DAY=One of the following:
\begin{itemize}
\item day1=relative to the forecast on a time scale of 1d in the future (timestep=10m)
\item day2=relative to the forecast on a time scale of 2d in the future (timestep=1h)
\item day3=relative to the forecast on a time scale of 3d in the future (timestep=3h)
\end{itemize}


EXAMPLE OUTPUTS (each day has 18 outputs -> 6 outputs x 3 days):
\begin{verbatim}
1. 20210111-00_20210111-2339_20210112-1000_detLONG_standardALL_day1.csv (7 figures)
2. 20210111-00_20210111-2339_20210112-1000_probLONG_pwv_day1.csv (1 figure)
3. 20210111-00_20210111-2339_20210112-1000_probLONG_rh_night1.csv  (1 figure)
4. 20210111-00_20210111-2339_20210112-1000_probLONG_wd_night1.csv. (1 figure)
5. 20210111-00_20210111-2339_20210112-1000_probLONG_ws_night1.csv. (1 figure)
6. 20210111-00_20210111-2339_20210112-1000_probdetLONG_cloud_night1.csv (2 figures)
7. 20210111-00_20210112-2339_20210113-1001_detLONG_standardALL_night2.csv
8. 20210111-00_20210112-2339_20210113-1001_probLONG_pwv_night2.csv
9. 20210111-00_20210112-2339_20210113-1001_probLONG_rh_night2.csv
10. 20210111-00_20210112-2339_20210113-1001_probLONG_wd_night2.csv
11. 20210111-00_20210112-2339_20210113-1001_probLONG_ws_night2.csv
12. 20210111-00_20210112-2339_20210113-1001_probdetLONG_cloud_night2.csv
13. 20210111-00_20210113-2339_20210114-1001_detLONG_standardALL_night3.csv
14. 20210111-00_20210113-2339_20210114-1001_probLONG_pwv_night3.csv
15. 20210111-00_20210113-2339_20210114-1001_probLONG_rh_night3.csv
16. 20210111-00_20210113-2339_20210114-1001_probLONG_wd_night3.csv
17. 20210111-00_20210113-2339_20210114-1001_probLONG_ws_night3.csv
18. 20210111-00_20210113-2339_20210114-1001_probdetLONG_cloud_night3.csv
\end{verbatim}

FILES HEADER

The headers are mostly auto-esplicative.

We add here some more infos:
\begin{itemize}
\item  The files reporting the probability that seeing and tau0 fall within specific categories are indicated with Category 1, 2, 3, 4, 5, 6 and 7 according to file readme\_fig4.png and according to the SOW (September update).
\item According to SOW, section 4.1.6, the cloud cover is treated differently than the other parameters. Files whose name is *probdet\_cloud* include the determinist and probabilistic forecast according to Section 4.1.6 of SOW and the digital document AD2.
\item According to SOW - section 4.2.6 we include the probability that the wind direction blows from N, S, W and E. Meteo Convention: 0 deg = North, 90 deg = East. In our analysis we divided the 360 degrees in 4 sectors: "North": means that the wind direction is included between -45 degrees and + 45 degrees. Same logic for the others cardinal directions.
\item  The ranges for the calculation of the probability that all the atmospheric and astroclimatic parameters assume specific values match tables reported in SOW.
\end{itemize}


%\begin{itemize}
%\item  The files reporting the probability that seeing and tau0 fall within specific categories are indicated with Category 1, 2, 3, 4, 5, 6 and 7 according to file readme_fig4.png and according to the SOW (September update).
%\item According to SOW, section 4.1.6, the cloud cover is treated differently than the other parameters. Files whose name is *probdet_cloud* include the determinist and probabilistic forecast according to Section 4.1.6 of SOW and the digital document AD2.
%\item According to SOW - section 4.2.6 we include the probability that the wind direction blows from N, S, W and E.  Meteo Convention: 0 deg = North 90 deg = East. In our analysis we divided the 360 degrees in 4 sectors: "North": means that the wind direction is included between -45 degrees and + 45 degrees. Same logic for the others cardinal directions. 
%\item The ranges for the calculation of the probability that all the atmospheric and astroclimatic parameters assume specific values match tables reported in SOW.
%\end{itemize}




%%%%%%%%%%%%%%%%%%%%%%%%%%%%%%%%%
\clearpage
%\begin{center}
%\HRule \\[0.4cm]
%\begin{figure}[htbp]
%\centering
%{\includegraphics[width=2.5cm]{/home/report/scripts/LOGOS/fate_logo_11def}}
%\caption{Add some meaningfull logo or figure here}
%\end{figure}
%\HRule \\[0.4cm]
%\end{center}

%\section{Contacts}
\vspace{1.5cm}
\begin{flushleft}
\HRule \\[0.4cm]
{\includegraphics[width=2.5cm]{/home/report/SCRIPTS/LOGOS/fate_logo_11def}}
\\
\Large{For any further information or questions please contact:
\\ WHO
\\ ADDRESS1
\\ ADDRESS2
\\ EMAIL: \href{mailto:e.musk@lamma.toscana.it}{e.musk@lamma.toscana.it}
\\ WEB: \href{www.lamma.toscana.com}{www.lamma.toscana.com}
\\ PHONE: MY\_PHONE
\\ FAX: MY\_FAX(???)
\\ OTHER: \dots}
\HRule \\[0.4cm]
\end{flushleft}

%%%%%%%%%%%%%%%%%%%%%%%%%%%%%%%%%

\end{document}

